% MMW-METHOD-CONTRACT-SHA256: d7b5c81661c8226b159e9d031c95d5cadec92a2b327791e0f6c2fb3b01dffc48
\documentclass[withoutpreface,bwprint]{cumcmthesis}

\usepackage{amsmath,amssymb}
\usepackage{graphicx}
\graphicspath{{figures/}{./}}
\usepackage{booktabs}
\usepackage{multirow}
\usepackage{longtable}
\usepackage{tabularx}
\usepackage{float}
\usepackage{subcaption}
\usepackage{url}
\usepackage{listings}
\hypersetup{hidelinks}

\title{面向算电协同的多目标调度优化研究}
\author{参赛队号：\quad 题号：C}
\date{}

\begin{document}

\maketitle

\begin{abstract}
本文研究六区域异构 AI 任务、可再生能源、储能与电网的协同调度。针对问题一，采用严格的训练（0--2351）、验证（2352--2375）和测试（2376--2399）时间切分，在季节朴素模型与 $3\times3$ Huber 参数网格间逐序列选型，并执行闭环递推；18 条序列的平均 MAE、RMSE 和 WAPE 分别为 25.5892~GPU、34.5070~GPU 和 0.8231。针对问题二，构造成本、碳、新能源、服务与平衡五类增量邻域，各评价 $10000$ 个候选并对最终 $50000$ 个任务执行全量约束审计；平衡方案净购售电成本为 $-421801727.45$~元，碳排放为 $16.5160$~tCO$_2$，新能源利用率为 $0.679962$。针对问题三，在含 no-action 的七个储能策略间用离散 SOC 动态规划搜索，经轨迹去重和支配检验选择平衡方案，其成本为 $-458584632.33$~元、碳排放为 $0.2241$~tCO$_2$。针对问题四，对 17 个情景逐一从共同基线重启，每轮根据上一轮储能轨迹重建 600 个任务候选，最多 8 轮、连续 2 轮无改善停止；共评价 $22200$ 个候选，基准情景实时 SLA 与按期完成率均为 $1$。两次隔离运行除实际耗时外 48 个正式输出文件逐字节一致，硬约束最大残差不超过 $10^{-6}$。所得结论属于有限搜索下的 scenario-feasible 结果，不声称全局最优。
\end{abstract}

\noindent\textbf{关键词：}异构 AI 任务；算力调度；Huber 回归；储能协同；碳感知调度；情景分析

\clearpage

\small
\linespread{0.94}\selectfont
\setlength{\parskip}{0pt}
\setlength{\emergencystretch}{3em}
\setlength{\abovedisplayskip}{4pt plus 1pt minus 1pt}
\setlength{\belowdisplayskip}{4pt plus 1pt minus 1pt}
\setlength{\abovedisplayshortskip}{3pt plus 1pt minus 1pt}
\setlength{\belowdisplayshortskip}{3pt plus 1pt minus 1pt}
\setlength{\textfloatsep}{8pt plus 2pt minus 2pt}
\setlength{\floatsep}{7pt plus 2pt minus 2pt}
\setlength{\intextsep}{7pt plus 2pt minus 2pt}
\setlength{\abovecaptionskip}{3pt}
\setlength{\belowcaptionskip}{2pt}

\section{问题重述}

随着人工智能推理与训练任务持续增长，数据中心的算力分配已经与能源供给、碳排放和通信服务质量形成紧密耦合。题目给出了六个区域的 GPU 容量、IT 与设施功率边界、逐时非 AI 负荷、PUE、电价、碳强度、新能源出力、购售电限制、储能参数以及区域间单向网络时延，同时提供了三类异构 AI 任务的到达时刻、GPU 需求、预计持续时间、来源区域、最大允许时延和完成期限。需要在满足物理边界与服务约束的条件下，合理确定任务的执行区域和开工时刻，并协调新能源、储能及电网购售电行为。

整个运行过程具有明确的分层时间边界。第 $0$ 至第 $2399$ 小时构成任务到达主时域，第 $2400$ 至第 $2405$ 小时不再接收新任务，仅允许此前到达的弹性任务继续执行并完成结清；第 $2406$ 小时不得安排任何任务，只用于电力运行和储能终端状态结算。所有任务均不可拆分、不可抢占，也不得在执行过程中迁移。实时推理任务必须在到达时刻立即开工，批量推理与 AI 训练任务可以在允许范围内延迟，但其完成时点不得超过任务截止时间和系统统一终点。

对于问题一，需要按区域和任务类型统计 GPU 需求的规模、分布及时间规律，建立短期需求预测模型，并按照时间顺序完成模型训练、验证与最终测试。预测结果仅用于模型精度评价，最后 $24$ 小时的基础调度必须使用第 $2376$ 至第 $2399$ 小时的实际到达任务。调度时还需根据分钟级持续时间计算任务与各小时的真实重叠，据此分别审计瞬时 GPU 占用和小时电量，并展示任务实际起止时刻、执行区域及区域 GPU 利用情况。

对于问题二，需要使用全体实际任务和逐时电力参数重新构造碳感知调度方案。在任务唯一执行、实时任务即时开工、网络时延、GPU 容量、IT 功率、设施功率及完成时限均得到满足的基础上，系统运行成本、购电碳排放、有效网络时延和新能源利用率应被统一复算。问题二不配置储能，新能源被用于直接供能、区域外送或弃置，不足负荷由电网购电补足，且购电与售电行为必须服从区域边界及互斥要求。

对于问题三，任务调度结果不再作为决策变量，区域 AI IT 负荷被固定为附件中的基准负荷，区域设施负荷由基准 AI IT 负荷与非 AI IT 负荷之和乘以 PUE 得到。需要在储能初始状态、荷电状态递推、容量和功率边界、充放电互斥、购售电互斥以及终端能量不透支等条件下，确定逐区域逐时的充电、放电、购电、售电、新能源分配和弃电策略，并将含储能方案与无储能比较状态在成本、碳排放、峰值净购电功率和负荷爬坡方面进行比较。

对于问题四，需要把任务迁移与开工、设施负荷、新能源分配、储能运行以及购售电决策纳入同一框架，在运行成本、碳排放、网络时延、服务质量、新能源利用率和区域峰值净购电之间进行权衡。碳约束、电价机制与新能源波动等情景应从同一联合基准独立重构，并在相同任务集、容量边界、功率映射和评价口径下接受完整审计，从而辨析不同外部条件对任务迁移方向、执行窗口与能源行为的影响。

\section{问题分析}
本题包含三个不同尺度：任务在离散区域和整数开工时刻上选择，分钟级持续时间决定逐小时重叠，能源与储能则按小时结算。直接把任务汇总为平均负荷会掩盖瞬时 GPU 和功率冲突，因此本文始终保留任务级轨迹，并由轨迹重建区域--小时资源数组。

问题一的关键是防止未来信息泄漏。所有滞后特征只使用预测时点之前的数据，验证集只负责模型选择，测试集不参与任何超参数决策。问题二的组合空间远大于可用计算预算，故采用确定性增量评价：一次移动只更新受影响的区域--小时数组，但每个命名方案结束后仍对全部任务重新审计。问题三在给定设施负荷后按区域分解，以有限 SOC 状态动态规划产生不同目标偏好的储能轨迹；no-action 必须进入候选集。问题四不是在固定三套任务方案中挑选，而是对每个情景从共同基线独立重启，交替执行储能优化和由当前储能边际信号驱动的任务邻域重建。
\section{模型假设}

为统一小时级调度与分钟级资源计量，假设任务开工时刻按照整数小时选取，而任务在各小时内的 GPU-hour 和能耗根据实际起止时刻与该小时的重叠时长计算。附件中的 GPU 需求、预计持续时间、到达时刻、来源区域、最大允许时延和最晚完成时刻均被视为准确且确定的输入参数，不另行建立执行时间和资源需求的随机模型。所有任务仅能选择一个执行区域和一个开工时刻，自开工至完成保持连续运行，不允许拆分、抢占、暂停或中途迁移。

假设实时推理任务必须在到达时刻立即开工，其执行区域只能从满足单向网络时延和资源容量的区域中选择。批量推理与 AI 训练任务可以延迟开工，但不得早于到达时刻，且必须在任务有效截止时点与系统时点 $2406$ 中较早者之前完成。第 $2400$ 至第 $2405$ 小时只用于此前到达弹性任务的结清，第 $2406$ 小时不产生任务资源占用，只承担电力与储能终端结算。

假设同一任务类型的单位等效 GPU 平均 IT 功率在研究期内保持为功率映射附件给定值，不随区域和时间改变。非 AI IT 负荷在各区域和各时段固定且不可迁移，只有 AI 任务负荷能够随执行区域及开工时刻改变。区域设施功率由 AI 与非 AI 合计 IT 功率乘以对应 PUE 得到，PUE 在研究期内按附件取值，不建立其随环境或负荷率变化的动态模型。

假设区域间通信只体现为附件给出的来源区域至执行区域单向网络时延，时延在研究期内确定且不受网络拥塞影响。来源区域与执行区域相同时仍采用时延矩阵中的对应记录，不自行设定为零。由于题面没有给出网络带宽、迁移数据量、传输能耗和传输费用等信息，上述因素不被纳入模型。

假设购电价格、售电价格、电网碳强度、可用新能源和固定负荷在单个小时内保持恒定。各区域逐小时独立满足能源平衡，不建立区域间线路潮流模型，区域售电被视为向系统边界外送。可用新能源可以被分配至直接供给负荷、储能充电、外送或弃电，其优先关系由实际策略与能源平衡共同确定，而不预设全部新能源均被本地直接消纳。运行碳排放仅按电网购电量与相应碳强度累计，不计新能源、储能设备和通信设施的生命周期排放。

假设储能在第 $0$ 小时运行前的状态唯一取自附件中的初始 SOC，后续状态均在该绝对值基础上递推，不将初始状态归零或作为待估参数。区域逐时基准 SOC 只用于方案核验与比较，不作为优化轨迹的固定约束。充电效率作用于进入储能的能量，放电效率作用于储能侧能量消耗，未给出的自放电与储能退化成本不予考虑。每个区域在第 $2406$ 小时运行后的 SOC 不低于其初始 SOC，从而避免通过透支终端能量获得虚假的当期收益。

假设同一储能在同一小时只能处于充电、放电或空闲状态，不能同时充放电；区域购电和售电也不能在同一小时同时发生。两类行为均受附件所给功率上限约束，若价格和效率条件不能自然排除无效循环，则通过互斥状态约束加以限制。问题三中的 IT 负荷固定为基准 AI IT 负荷与非 AI IT 负荷之和，不重新调整任务的执行区域和开工时刻。

假设预测模型与正式调度严格分离，预测区间的估计值只用于测试精度评价，最后 $24$ 小时基础调度以及问题二至问题四的正式计算均使用实际任务和实际逐时参数。服务质量仅由题内可观测的即时开工、单向时延 SLA、按期完成情况及弹性任务等待时间表示，不引入用户满意度或可靠性等缺乏数据支持的变量。进行碳约束、电价机制与新能源波动情景比较时，除指定变化因素外，其余任务集、设备容量、网络时延、功率映射及评价公式均保持一致。

假设普通空白值在未经核验前均属于不可裁决缺失，不自动替换为零。只有在互斥进流与出流字段中，空白所在行为孤立事件、对应另一列明确非零，且备注、相邻记录和全表口径共同支持互斥解释时，才将其认定为结构零。经过处理的数据还应接受任务唯一执行、实时即时开工、时延与截止、GPU 和功率容量、能源守恒、SOC 递推、终端状态及购售电边界的统一审计。

\section{符号说明}

本文使用的主要符号如表~\ref{tab:symbols} 所示。对于同时具有输入参数和派生量的符号，其具体属性在表中一并注明；功率变量均以区域、小时为索引，能量累计时采用时段长度 $\Delta t$ 进行换算。

{\scriptsize
\renewcommand{\arraystretch}{0.88}
\setlength{\tabcolsep}{3pt}
\begin{longtable}{lll}
\caption{主要符号及其含义}
\label{tab:symbols}\\
\toprule
符号 & 含义 & 类型或单位 \\
\midrule
\endfirsthead
\toprule
符号 & 含义 & 类型或单位 \\
\midrule
\endhead
\bottomrule
\endfoot
$i,r,t,s,k$ & 任务、区域、小时、开工时刻和任务类型索引 & 索引 \\
$\mathcal R$ & 六个计算区域构成的集合 & 集合 \\
$\mathcal K$ & 三类异构 AI 任务构成的集合 & 集合 \\
$\mathcal T_A$ & 任务到达主时域 & $\{0,\ldots,2399\}$ \\
$\mathcal T_C$ & 末端任务结清时域 & $\{2400,\ldots,2405\}$ \\
$\mathcal T_E$ & 电力与储能结算时域 & $\{0,\ldots,2406\}$ \\
$\mathcal S$ & 合法整数开工时刻集合 & $\{0,\ldots,2405\}$ \\
$\Delta t$ & 单个电力结算时段长度 & h \\
$a_i$ & 任务 $i$ 的到达小时 & h \\
$d_i,h_i$ & 任务持续分钟数及换算后的持续小时数 & min，h \\
$g_i$ & 任务 $i$ 持续占用的等效 GPU 数 & GPU \\
$k_i,o_i$ & 任务 $i$ 的类型和来源区域 & 类别，区域 \\
$F_i$ & 任务 $i$ 的有效完成截止时点 & h \\
$W_i^{\max}$ & 任务 $i$ 的最大允许等待时间 & h \\
$\ell_i^{\max}$ & 任务 $i$ 允许的最大网络时延 & ms \\
$\ell_{or}$ & 来源区域 $o$ 至执行区域 $r$ 的单向时延 & ms \\
$\omega_{it}(s)$ & 任务 $i$ 在开工时刻 $s$ 下与小时 $t$ 的重叠时长 & h \\
$\chi_{it}(s)$ & 任务 $i$ 在小时 $t$ 是否存在瞬时占用 & 二元参数 \\
$\Omega_i$ & 任务 $i$ 的合法区域--开工候选域 & 有限集合 \\
$x_{irs}$ & 任务 $i$ 是否在区域 $r$ 于时刻 $s$ 开工 & 二元变量 \\
$s_i,c_i,r_i$ & 任务实际开工时刻、完成时刻和执行区域 & h，h，区域 \\
$D_{rkt}$ & 区域 $r$ 类型 $k$ 在小时 $t$ 的实际到达 GPU 需求 & GPU \\
$D'_{rkt}$ & 用于滞后回归的历史需求输入 & GPU \\
$\widehat D_{rkt}$ & 区域--任务类型逐小时 GPU 需求预测值 & GPU \\
$u_{rkt}$ & Huber 回归截断前的线性预测量 & GPU \\
$z_t$ & 小时 $t$ 的日周期与周周期特征向量 & 无量纲 \\
$\beta,\gamma,\alpha_r,\delta_k$ & 回归系数、周期系数、区域效应和类型效应 & 参数 \\
$\lambda$ & Huber 回归的正则化系数 & 无量纲 \\
$\sigma_{rk}$ & 区域--任务类型序列的残差尺度 & GPU \\
$\rho_{\delta}(\cdot)$ & 阈值为 $\delta$ 的 Huber 损失函数 & 无量纲 \\
$G_r^{\max}$ & 区域 $r$ 的可调度 GPU 容量 & GPU \\
$P_r^{IT,\max}$ & 区域 $r$ 的最大 IT 功率 & MW \\
$P_r^{F,\max}$ & 区域 $r$ 的最大设施功率 & MW \\
$p_k^{GPU}$ & 类型 $k$ 的单位等效 GPU 平均 IT 功率 & MW/GPU \\
$\pi_r$ & 区域 $r$ 的能源使用效率 PUE & 无量纲 \\
$L_{rt}^{N}$ & 区域 $r$ 的不可迁移非 AI IT 负荷 & MW \\
$L_{rt}^{BAI}$ & 问题三固定的基准 AI IT 负荷 & MW \\
$L_{rt}^{AI,cap}$ & AI 任务形成的瞬时 IT 功率 & MW \\
$L_{rt}^{AI,avg}$ & AI 任务形成的小时平均 IT 负荷 & MW \\
$E_{rt}^{AI}$ & AI 任务在小时 $t$ 消耗的电量 & MWh \\
$L_{rt}^{F,avg}$ & 区域设施侧小时平均负荷 & MW \\
$U_{rt}^{peak}$ & 区域小时峰值 GPU 利用率 & 无量纲 \\
$U_{rt}^{avg}$ & 区域小时 GPU-hour 平均利用率 & 无量纲 \\
$R_{rt}^{av}$ & 区域可用新能源功率 & MW \\
$R_{rt}^{use}$ & 直接满足设施负荷的新能源功率 & MW \\
$R_{rt}^{ch}$ & 用于储能充电的新能源功率 & MW \\
$R_{rt}^{sell}$ & 新能源外送功率 & MW \\
$R_{rt}^{curt}$ & 弃置的新能源功率 & MW \\
$G_{rt}^{load}$ & 电网直接供给设施负荷的功率 & MW \\
$G_{rt}^{ch}$ & 电网供给储能充电的功率 & MW \\
$G_{rt}^{buy}$ & 区域从电网购入的总功率 & MW \\
$G_{rt}^{sell}$ & 区域向系统边界外送的功率 & MW \\
$I_r^{\max}$ & 区域最大购电功率 & MW \\
$X_r^{grid}$ & 区域电网侧最大外送功率 & MW \\
$X_r^{storage}$ & 储能信息限定的售电功率上限 & MW \\
$y_{rt}^{grid}$ & 区域购电与售电互斥状态 & 二元变量 \\
$C_{rt}$ & 储能总充电功率 & MW \\
$D_{rt}$ & 储能放电功率 & MW \\
$S_{rt}$ & 区域储能在小时 $t$ 运行后的 SOC & MWh \\
$B_r$ & 区域储能额定容量 & MWh \\
$S_r^0$ & 第 $0$ 小时运行前的初始 SOC & MWh \\
$S_r^{min},S_r^{max}$ & 区域储能 SOC 下限与上限 & MWh \\
$C_r^{\max},D_r^{\max}$ & 区域最大充电与放电功率 & MW \\
$\eta_r^c,\eta_r^d$ & 区域储能充电效率与放电效率 & 无量纲 \\
$N_{rt}$ & 区域净购电功率 $G_{rt}^{buy}-G_{rt}^{sell}$ & MW \\
$P_r^{peak}$ & 区域时域内非负峰值净购电功率 & MW \\
$V_r$ & 区域净购电绝对爬坡量 & MW \\
$C_{\mathrm{op}}$ & 购电支出扣除售电收入后的运行成本 & 元 \\
$E_{\mathrm{CO2}}$ & 电网购电产生的累计碳排放 & tCO$_2$ \\
$L_{\mathrm{net}}$ & 按 GPU 需求与执行时长加权的网络时延 & ms \\
$W_{\mathrm{net}}$ & 网络时延指标的有效权重和 & GPU$\cdot$h \\
$U_{\mathrm{RE}}$ & 新能源利用率 & 无量纲 \\
$W_{\mathrm{RE}}$ & 新能源利用率的可用新能源总量 & MWh \\
$Q_{\mathrm{RT}}$ & 实时任务即时开工满足率 & 无量纲 \\
$Q_{\mathrm{SLA}}$ & 网络时延 SLA 满足率 & 无量纲 \\
$Q_{\mathrm{deadline}}$ & 任务按期完成率 & 无量纲 \\
$Q_{\mathrm{wait}}$ & 弹性任务等待质量指标 & 无量纲 \\
$Q_{\mathrm{service}}$ & 综合服务质量指标 & 无量纲 \\
$\Psi_m$ & 实际有限候选搜索使用的代理综合评分 & 无量纲 \\
$w_j$ & 有效评价维度的实际等权系数 & 无量纲 \\
$d_j$ & 实际尺度 $\max\{|a_{j,m_0}|,1\}$ & 对应指标单位 \\
$a_{jm}$ & 候选方案 $m$ 在指标 $j$ 上的取值 & 对应指标单位 \\
$m_0$ & 标准化评分使用的基准方案索引 & 索引 \\
$N_{\mathrm{back}}$ & 冲突回溯节点预算 & 节点 \\
$N_{\mathrm{move}}^{q2}$ & 问题二单任务候选评价预算 & 次 \\
$N_{\mathrm{move}}^{q3}$ & 问题三储能动作候选评价预算 & 次 \\
$N_{\mathrm{move,total}}^{q4}$ & 问题四基准与全部情景共享的评价预算 & 次 \\
$D_{\mathrm{search}}$ & 启发式搜索的内部截止时刻 & 时间点 \\
$\rho_C$ & 情景中的碳约束水平参数 & 无量纲 \\
$\delta_P$ & 峰谷电价差调整系数 & 无量纲 \\
$\delta_S$ & 售电价格调整系数 & 无量纲 \\
$\delta_R$ & 新能源出力水平调整系数 & 无量纲 \\
$\xi$ & 新能源波动情景中的固定种子随机扰动 & 无量纲 \\
$\kappa_{rt}$ & 区域逐时电网购电碳强度 & tCO$_2$/MWh \\
$c_{rt}^{buy},c_{rt}^{sell}$ & 区域逐时购电价格和售电价格 & 元/MWh \\
$\mathrm{EFC}_r$ & 区域储能的等效完整循环次数 & 次 \\
\end{longtable}
}

表~\ref{tab:symbols} 同时区分了瞬时容量变量与小时能量变量，其中 $\chi_{it}(s)$ 被用于 GPU、IT 功率和设施功率的峰值审计，$\omega_{it}(s)$ 则被用于 GPU-hour 与 AI 电量计算。该区分能够保留非整小时任务的真实资源占用，避免小时平均功率掩盖瞬时超限。

能源变量按来源和去向被进一步拆分，使新能源守恒、电网购电成本和储能充电来源能够分别核算。由此，新能源利用率不会把弃电计入有效利用量，碳排放也不会遗漏电网为储能充电所提供的电量。

评分符号仅作用于分母有效且在候选之间存在真实差异的指标。若 $W_{\mathrm{net}}$、$W_{\mathrm{RE}}$ 或服务质量分量的相应分母不为正，则该指标被标记为不可用并退出综合评分，从而避免无依据的除零处理。

% MMW-METHOD-CONTRACT: 17e0e931dd01e167c81738be87e8eca8e00933987071026ea2be151827f2431e
% MMW-METHOD-CONTRACT: ead1dc31de5453e83924fd80c141e14521697c804489a275435dbb5338b47dac
% MMW-ID: CON-Q1-1
% MMW-ID: CON-Q1-2
% MMW-ID: CON-Q1-3
% MMW-ID: CON-Q1-4
% MMW-ID: CON-Q1-5
% MMW-ID: CON-Q1-6
% MMW-ID: CON-Q1-7
% MMW-ID: CON-Q2-1
% MMW-ID: CON-Q2-2
% MMW-ID: CON-Q2-3
% MMW-ID: CON-Q2-4
% MMW-ID: CON-Q2-5
% MMW-ID: CON-Q2-6
% MMW-ID: CON-Q3-1
% MMW-ID: CON-Q3-2
% MMW-ID: CON-Q3-3
% MMW-ID: CON-Q3-4
% MMW-ID: CON-Q3-5
% MMW-ID: CON-Q3-6
% MMW-ID: CON-Q3-7
% MMW-ID: CON-Q4-1
% MMW-ID: CON-Q4-2
% MMW-ID: CON-Q4-3
% MMW-ID: CON-Q4-4
% MMW-ID: CON-Q4-5
% MMW-ID: CON-Q4-6
% MMW-ID: OBJ-Q1
% MMW-ID: OBJ-Q2
% MMW-ID: OBJ-Q3
% MMW-ID: OBJ-Q4
% MMW-ALGORITHM: 严格闭环预测；五方案增量任务局部搜索；离散SOC动态规划；按情景动态重建邻域的任务-储能交替有限搜索
\section{模型的建立与求解}
\subsection{任务与能源统一约束}
令 $x_i=(r_i,s_i)$ 表示任务 $i$ 的执行区域与开工时刻，$f_i=s_i+d_i$。每个任务只执行一次；实时任务满足 $s_i=a_i$，弹性任务满足 $a_i\le s_i$ 且 $f_i\le\min(l_i,2406)$。区域间实际时延不得超过任务上限。对区域 $r$、小时 $t$，以区间重叠指示量累计 GPU 与瞬时 IT 功率，并要求
\[
G_{rt}\le \bar G_r,\qquad L^{\rm nonAI}_{rt}+L^{\rm AI}_{rt}\le \bar P^{\rm IT}_r,
\]
且乘以 PUE 后不超过设施功率边界。能源侧满足逐时守恒，购电、售电、充电、放电均非负，且购售电与充放电分别互斥。储能状态满足
\[
E_{r,t+1}=E_{rt}+\eta_r^c P^c_{rt}-P^d_{rt}/\eta_r^d,
\quad E_{r,2407}=E_{r0}.
\]
成本按购电支出减售电收入计算，碳排放只累计购电量与碳强度的乘积；负成本表示题面核算边界内售电收入高于购电支出，并不表示物理能耗为负。

\subsection{问题一：闭环预测}
对六区域三任务类型形成 18 条逐时 GPU 到达需求序列。候选包括日周期与周周期季节朴素模型，以及 $\epsilon\in\{1.10,1.35,1.75\}$、$\alpha\in\{10^{-4},10^{-3},10^{-2}\}$ 的 Huber 稳健模型\cite{huber1964robust}。训练、验证、测试边界分别为 2351、2352--2375 和 2376--2399；多步预测逐点回填预测值而非真实未来值。最终 17 条序列选择 Huber，1 条选择 168 小时季节朴素模型，平均 MAE=25.5892、RMSE=34.5070、WAPE=0.8231。图~\ref{fig:gpu} 展示最终任务方案的区域 GPU 利用率。
\begin{figure}[H]\centering
\includegraphics[width=.86\textwidth]{fig_1_gpu_utilization.png}
\caption{区域 GPU 利用率（完整任务轨迹复算）}\label{fig:gpu}
\end{figure}

\subsection{问题二：五方向有限邻域}
对成本、碳、新能源、服务和平衡五个方向分别检查 $10000$ 个任务移动候选。候选通过增量数组判断容量与能源变化，结束后再对 50000 个任务执行唯一性、时窗、时延、GPU、IT、设施功率和能源边界全量审计。表~\ref{tab:q2} 给出结果；方案名表示搜索方向，不表示跨方案或完整空间中的对应目标最优值。
\begin{table}[H]\centering\small
\caption{问题二五类搜索结果}\label{tab:q2}
\begin{tabular}{lrrrrr}\toprule
方向 & 成本/亿元 & 碳/t & 新能源率 & 时延/ms & 接受移动\\\midrule
成本 & -4.2334 & 10.2851 & 0.680386 & 6.2235 & 422\\
碳 & -4.1966 & 14.3916 & 0.679078 & 5.1594 & 0\\
新能源 & -4.2230 & 14.3916 & 0.680193 & 7.0758 & 761\\
服务 & -4.1966 & 14.3916 & 0.679078 & 5.1594 & 0\\
平衡 & -4.2180 & 16.5160 & 0.679962 & 6.5824 & 854\\\bottomrule
\end{tabular}
\end{table}
平衡方案迁移 561 个任务，服务质量为 0.990393。五次搜索合计恰好评价 50000 个候选，每个方向均单独通过一次全量任务审计。

\subsection{问题三：离散 SOC 动态规划与 Pareto 筛选}
每个区域使用 21 个等距 SOC 状态并显式加入初始 SOC，逐时枚举可达状态转移。候选包括 no-action、成本、碳、新能源、峰值、爬坡和平衡七类策略；相同轨迹仅计一次，并按多目标支配关系筛选\cite{miettinen1999nonlinear}。七个候选形成 6 条唯一轨迹，平衡方案通过候选集内非支配检验，成本、碳、新能源利用率、峰值和爬坡分别为 $-458584632.33$ 元、$0.224084$ tCO$_2$、0.771439、0.373986 MW 和 0.747972 MW。no-action 对应成本 $-419820327.74$ 元、碳排放 8.850563 tCO$_2$。图~\ref{fig:soc} 给出选中轨迹。
\begin{figure}[H]\centering
\includegraphics[width=.86\textwidth]{fig_2_storage_soc.png}
\caption{问题三选中储能轨迹的 SOC}\label{fig:soc}
\end{figure}

\subsection{问题四：动态任务--储能交替}
17 个情景包括基准、四个碳约束水平、三组峰谷价差、三组售电机制、三组新能源水平及种子 2026、2027、2028 的新能源波动。每个情景从相同任务基线和真实初始 SOC 独立重启；每轮先求储能，再把电价、碳强度、购电与弃电组合为任务移动信号，重新生成并评价 600 个任务候选。最多 8 轮，连续 2 轮无改善停止。

基准情景完成 2 轮、1200 次候选评价，成本 $-458341900$ 元、碳排放 0.578894 tCO$_2$、新能源利用率 0.768912、峰值购电 0.443765 MW；实时 SLA 和按期完成率均为 1，最大等待 6 h。10\% 碳削减情景在第 2 轮接受新任务方案，碳排放降至 0.424511 tCO$_2$；20\% 与 30\% 目标在声明启发式内未达到，不能据此判定数学不可行。全部情景合计 37 轮、22200 次候选评价，且每轮的候选与 incumbent 哈希均被记录。



\subsection{方法偏差与适用边界}
% MMW-LIMITATION: D1
\paragraph{方法边界 D1}Q2/Q4为预声明候选预算内最佳可行解，不是穷举最优解

% MMW-LIMITATION: D2
\paragraph{方法边界 D2}Q3 SOC使用有限状态离散，未计入题面未提供的退化成本和循环寿命硬上限

% MMW-LIMITATION: D3
\paragraph{方法边界 D3}Q4随机情景另使用固定种子2027和2028并逐情景独立重启

% MMW-LIMITATION: L1
\paragraph{方法边界 L1}无独立Oracle，最高可信等级为scenario-feasible

% MMW-LIMITATION: L2
\paragraph{方法边界 L2}储能SOC离散化且缺少退化成本/循环寿命数据

% MMW-LIMITATION: L3
\paragraph{方法边界 L3}Q2各命名方案表示搜索方向，不表示跨方案的对应目标最小值

\section{模型的检验与灵敏度分析}
\subsection{硬约束与双运行检验}
任务审计最大 GPU、唯一性、时延、截止、实时开工和终端残差均不超过 $10^{-6}$；Q3 与 Q4 储能能量平衡、SOC、功率、互斥和终端状态的最大绝对残差均为 $2.274\times10^{-13}$。两次隔离运行使用相同输入副本和 SHA-256 为 \texttt{846f527b...} 的冻结程序，自然完成耗时分别为 1687.58 s 和 1706.49 s。除 method\_runtime.json 中的实际耗时外，48 个正式输出文件逐字节一致，results.json 的 SHA-256 为 \texttt{691e95d6...}。

\subsection{情景响应}
新能源水平为 0.8 时，碳排放升至 500.5503 tCO$_2$、峰值购电 92.1180 MW；水平为 1.2 时，两者均降为 0。售电价格系数从 0、0.5 到 1 时，净成本从 524.79 元降至 $-2.291688\times10^8$ 元和 $-4.583419\times10^8$ 元，说明结果对题面售电收益口径高度敏感。图~\ref{fig:ren} 给出新能源水平灵敏度。
\begin{figure}[H]\centering
\includegraphics[width=.82\textwidth]{sensitivity_renewable_level.png}
\caption{新能源水平灵敏度}\label{fig:ren}
\end{figure}

\section{模型评价}
本文的优点是：预测选型与测试严格隔离；任务侧使用增量计算但不省略最终 5 万任务全量审计；储能候选显式包含 no-action；Q4 每轮从上一轮储能结果动态重建邻域；全部结果具有结构化运行证据和两次独立复算哈希。

局限性也必须明确。第一，Q2/Q4 只在预声明有限候选预算内寻找 best-found 可行解，不具有全局最优证书；Q3 只在离散 SOC 网格和七类目标策略中做候选集 Pareto 判断。第二，题面未给储能退化成本与循环寿命上限，平衡轨迹的等效循环次数较高，不能直接作为真实设备控制指令。第三，通信模型只含固定时延，未计带宽拥塞、迁移数据量和传输能耗。第四，没有独立 Oracle 或真实部署数据，因此可信等级最高为 scenario-feasible。

\appendix
\section{程序代码与运行说明}
正式程序为交付目录中的 \texttt{solution.py}，在含 \texttt{附件数据} 的目录执行：
\begin{verbatim}
python solution.py
python validate_results.py
\end{verbatim}
程序无默认墙钟终止；只有显式设置正数环境变量 \texttt{MMW\_MAX\_RUNTIME\_SECONDS} 才允许外部时间上限，触发时结果必须标记 incomplete 并阻断审批。正式结果写入 \texttt{output/data}，图表写入 \texttt{output/figures}。主要证据包括 \texttt{results.json}、\texttt{method\_runtime.json}、五方向 Q2 比较表、Q3 轨迹表、Q4 情景表和逐轮交替 trace。

\bibliographystyle{plain}
\bibliography{references}

\end{document}

